
\textbf{Corollaire 6.14. (I. Hermann).} \emph{Pour tout $N\ge0$, il existe un polyn\^ome $R_N(m)$ tel que pour tout corps $k_0$ et tout id\'eal de l'anneau de polyn\^omes $k_0[T_1,\ldots,T_N]$, engendr\'e par des \'el\'ements de degr\'e $\le m$, la racine $\sqrt I$ de $I$ est engendr\'e par des \'el\'ements de degr\'e $\le R_N(m)$.}

En effet, soit $k$ la cl\^oture alg\'ebrique de $k_0$. Introduisant une ind\'etermin\'ee auxilaire, consid\'erons $Y$ (resp. $X$) le sous-sch\'ema de $P=\mathbb P^N_k$ d\'efini par l'id\'eal homog\'en\'eis\'e de $I$ (resp. de $\sqrt I$). Dans le polyn\^ome de Hilbert $\chi({O}_Y(n))=\sum_{i=0}^{N}b_i\binom{n+i}{i}$, on a $|b_i|\le c_i$ avec $c_i=A_{N-i}(0,\ldots,0,1;m)$ d'apr\`es 6.13.

Soit $X=\bigcup X^q$ la d\'ecomposition suivant la dimension. Alors $\text{degr\'e}(X^q)\le e_q$ avec $e_q=P_{N-q}(c_q,\ldots,c_N)$. En effet, quitte \`a couper $X$ (resp. $Y$) par un espace lin\'eaire ``g\'en\'eral'' de codimension $q$, on peut supposer $q=0$. Alors \'evidemment, $\text{degr\'e}(X^q)=h^0({O}_{X^q})$ et $h^0({O}_{X^q})\le h^0({O}_Y)$, et en vertu de 6.4 (iii), $h^0({O}_Y)\le e_q$.

Par cons\'equent, ${O}_{X^q}$ est un $(0,\ldots,0,e_q)$-faisceau gr\^ace \`a 6.10. Il r\'esulte donc de l'injection ${O}_X\to \prod{O}_{X^q}$ que ${O}_X$ est un $(e_0,\ldots,e_N)$-faisceau. Par suite, dans le polyn\^ome de Hilbert $\chi({O}_X(n))=\sum_{i=0}^{n}a_i\binom{n+i}{i}$, on a $|a_i|\le {}_i$ avec $f_i=A_{N-i}(e_i,\ldots,e_N;0)$ d'apr\`es 6.7. Enfin l'id\'eal $\sqrt I$ de $X$ est $R_N(m)$-r\'egulier avec $R_N(m)=P_N(f_0,\ldots,f_N)$ ; d'o\`u l'assertion gr\^ace \`a 1.3.

